\chapter{Research and Existing Solutions}
\label{ch:research}
\section{Research Method / Approach}
Research combined three activities: (i) reviewing the foundational literature on object detection, multi-object tracking, person re-identification, and planar visual metrology; (ii) surveying publicly described retail-analytics products and their stated capabilities; and (iii) a feasibility triage of candidate features against hardware (no dedicated GPU), timeline, and dependency constraints. Tools and libraries were then selected by prototyping on the supplied store videos.

\section{Existing Solutions / Products}
\textbf{Commercial retail analytics.} Vendors such as RetailNext and Sensormatic (ShopperTrak) market in-store traffic counting, conversion, and shopper-behaviour analytics. These products generally rely on purpose-installed sensors or cameras and on vendor subscriptions. \textbf{Checkout-free and AI-camera retail} systems, such as Amazon's Just Walk Out technology, track shoppers and products throughout the store using dense camera and sensor installations. \textbf{Generic video analytics.} Camera vendors and open-source tooling offer people counting and zone occupancy, but typically stop at counts and do not link customers to staff interactions.

\textbf{Research building blocks.} The YOLO family \cite{redmon2016yolo} provides fast single-stage detectors with CPU-feasible small variants, available through the Ultralytics implementation \cite{ultralytics}. SORT \cite{bewley2016sort} and ByteTrack \cite{zhang2022bytetrack} associate detections across frames into tracks. Person re-identification across cameras is surveyed in \cite{ye2022reid} and typically depends on learned embeddings and substantial compute.

\section{Comparison of Existing Solutions}
\begin{table}[H]
\centering
\caption{Qualitative comparison of approaches (based on publicly described capabilities)}
\begin{tabular}{P{0.24\textwidth}P{0.17\textwidth}P{0.17\textwidth}P{0.17\textwidth}P{0.15\textwidth}}
\toprule
Criterion & Commercial analytics & Checkout-free systems & Generic people counters & RetailPulse\\
\midrule
Reuses existing CCTV & Partly & No & Often & Yes\\
Needs GPU or special hardware & Often & Yes & Varies & No (CPU)\\
Staff--customer interaction metrics & Limited & Not the focus & No & Yes\\
Zone dwell and heatmaps & Yes & Yes & Partly & Yes\\
POS integration & Yes & Integrated & No & File import\\
Cost to small retailer & High & Very high & Low & Low\\
Biometric identification & Varies & Varies & No & No\\
\bottomrule
\end{tabular}
\end{table}
\todo{verify vendor capability claims against current product documentation before final submission}

\section{Limitations and Gaps}
Existing options either require new capital investment or stop at counting. None of the reviewed approaches is aimed at a small retailer that wants to turn existing cameras into engagement analytics (approach time, unattended customers, staff idle time) on CPU-only hardware, while also joining those findings with POS data and avoiding biometric identification. Research-grade re-identification is accurate but compute-heavy, which does not suit the hardware constraint.

\section{Proposed Solution}
RetailPulse is a LangGraph-orchestrated multi-agent pipeline. Each agent owns one responsibility: detection, classification, engagement, trajectory, narration. An orchestrator reconciles all outputs against the same tracked entities so that the dashboard shows one consistent view. Lightweight YOLO detection with local ByteTrack identifiers keeps the pipeline CPU-friendly. A POS module adds commercial context. A conservative camera-handoff stage links cameras using direction, timing, and coarse appearance, and reports uncertainty explicitly.

\section{Complex Computing Problem}
The complex computing problem is to derive reliable, per-person service and movement analytics from low-cost, uncontrolled CCTV video under tight compute limits. The problem involves several interacting difficulties:
\begin{itemize}
\item Conflicting requirements: accuracy against CPU-only throughput, and informativeness against privacy (no biometrics).
\item Ambiguous ground truth: whether a staff member near a customer constitutes an ``interaction'' must be defined operationally from distance, orientation, and duration.
\item Consistency across components: the outputs of several agents must be reconciled per entity despite identity switches, occlusions, and disagreement between agents.
\item Cross-camera association with limited evidence: without biometric identifiers or a verified shared cross-camera floor map, and with timing uncertainty, associations must preserve uncertainty. The separate single-camera floor homography does not establish cross-camera geometric alignment.
\item Non-technical users: results must be presented in a way that managers can act on without over-trusting heuristics.
\end{itemize}

\section{Technical Challenges}
\begin{itemize}
\item \textbf{CPU-only inference:} frame sampling (for example, every fifth frame in the handoff pipeline) and small detector variants are needed to keep batch processing practical.
\item \textbf{Staff/customer separation without biometrics:} both learned appearance-based role detection and earlier zone rules require evaluation against clothing similarity and staff outside their expected zones.
\item \textbf{Perspective and metric scale:} pixel separations depend on image location; a calibrated floor mapping is required before estimating metres.
\item \textbf{Floor-contact uncertainty:} seated people, hidden feet, chairs, overlap and frame truncation can make a box-bottom estimate unreliable.
\item \textbf{Camera movement:} a reference homography can become misaligned. Insufficient static floor features mean unverified alignment, not evidence that the camera is stationary.
\item \textbf{Identity switches and occlusion:} ID switching in crowded scenes affects interaction durations and dwell times.
\item \textbf{Camera handoff:} direction, time-window filtering, and HSV/shape descriptors give weak evidence; lighting differences between cameras degrade colour similarity.
\item \textbf{Agent reconciliation:} the orchestrator must resolve conflicts between agents (for example, an entity labelled staff by one rule and customer by another).
\item \textbf{Security of the web tier:} manager-only POS data must not be unlocked by client-side state alone.
\end{itemize}

\section{Research Findings and Design Decisions}
\subsection{Geometric Basis of Floor-Distance Estimation}
Projective geometry provides a mapping between two views of a plane through a homography \cite{hartley2004multiple,opencvHomography}. Single-view metrology establishes broader foundations for recovering geometric measurements from perspective images using scene information \cite{criminisi2000metrology}; the implemented prototype specifically uses explicit image/floor landmark correspondences rather than reproducing that paper's complete procedure. RetailPulse selects planar floor mapping because the intended quantity is separation along the floor and the existing camera and detector can be reused. It does not estimate a depth map or reconstruct arbitrary 3D scene geometry.

Robust fitting with RANSAC \cite{fischler1981ransac} supports additional calibration correspondences and reference-alignment checks. ByteTrack \cite{zhang2022bytetrack} supplies per-camera association, while ORB \cite{rublee2011orb} supplies image features for the conservative motion monitor. These are established building blocks; the project's contribution is their integration with explicit position-quality states, configurable calibration and diagnostic exports. Published benchmark results for these methods are not claimed as RetailPulse performance.

\begin{longtable}{P{0.38\textwidth}P{0.5\textwidth}}
\caption{Findings and the decisions they influenced}\\
\toprule
Finding & Decision\\
\midrule
\endhead
YOLO detectors have small variants usable without a GPU \cite{redmon2016yolo,ultralytics}. & Use a lightweight YOLO-based detector and sample frames.\\
ByteTrack \cite{zhang2022bytetrack} gives stable local IDs without a separate appearance network. & Use ByteTrack for per-camera tracking.\\
Learned re-identification \cite{ye2022reid} is compute-heavy and biometric-adjacent. & Descope full re-identification. Ship a conservative two-camera handoff using direction, time window, and HSV/shape descriptors.\\
Single agents entangle concerns and are hard to test. & Split the pipeline into specialised agents and orchestrate with LangGraph.\\
Appearance evidence is not proof. & Report confirmed, uncertain, and unmatched associations. Do not force matches.\\
Client-side role flags are easily altered. & Issue server-signed HttpOnly session cookies and require an internal secret on direct backend calls.\\
Managers need explanations, not only numbers. & Add a narration agent and cited POS assistant answers.\\
Perspective changes pixel scale across the image \cite{hartley2004multiple,opencvHomography}. & Project supported foot proxies into floor metres before calculating separations.\\
Supplied floor dimensions determine metric scale. & Store editable landmarks and dimensions, with estimated inputs and approximate outputs identified in configuration.\\
Camera movement can invalidate the reference mapping. & Check static floor alignment using ORB/RANSAC; suppress output when shifted or unverified.\\
\bottomrule
\end{longtable}
